\section[Proof of Theorem 1: Intrinsic Attribute Preservation and Gradient Bounds]{Proof of Theorem 1: Intrinsic Attribute Preservation and Non-Vanishing Gradient Bounds in Heterogeneous Graph Neural Networks}
\label{proof:theorem1}

\begin{theorem}[Intrinsic Attribute Preservation and Gradient Lower Bounds]
Let $G = (V, E, \mathcal{T}_V, \mathcal{T}_E)$ be a directed heterogeneous multigraph with node type mapping $\tau: V \to \mathcal{T}_V$. Let each node $v \in V$ have an initial feature vector $x_v = h_v^{(0)} \in \mathbb{R}^{d_0}$. Consider an $L$-layer heterogeneous message-passing neural network where layer $l \in \{1, \dots, L\}$ computes hidden representations:
\begin{enumerate}
    \item \textbf{Intrinsic Feature Erasure in Relational Neighborhood Aggregation:}
    Under pure relational neighborhood aggregation without self-loops or skip connections:
    \begin{equation}
    h_v^{(l)} = \sigma \left( \sum_{r \in \mathcal{R}_{\text{in}}(\tau(v))} \bigoplus_{u \in \mathcal{N}_r(v)} \alpha_{vu}^{(l)} W_r^{(l)} h_u^{(l-1)} \right)
    \end{equation}
    where $\mathcal{N}_r(v) = \{u \in V \mid (u, r, v) \in E\}$ denotes incoming neighbors under relation $r$, the layer-1 hidden state $h_v^{(1)}$ is conditionally independent of $x_v$ given $\mathcal{N}(v)$. Consequently:
    \begin{equation}
    \frac{\partial h_v^{(1)}}{\partial x_v} = \mathbf{0} \in \mathbb{R}^{d_1 \times d_0}
    \end{equation}
    Moreover, for any source-only node $v$ satisfying $\mathcal{N}_{\text{in}}(v) = \emptyset$, $h_v^{(l)} = \mathbf{0}$ for all $l \ge 1$.
    \item \textbf{Guaranteed Gradient Lower Bound via Parameterized Residual Skips:}
    When parameterized residual skip connections are introduced:
    \begin{equation}
    \tilde{h}_v^{(l)} = h_v^{(l)} + W_{\text{skip}}^{(l)} \tilde{h}_v^{(l-1)}
    \end{equation}
    where $W_{\text{skip}}^{(l)} \in \mathbb{R}^{d_l \times d_{l-1}}$ has minimum singular value $\sigma_{\min}(W_{\text{skip}}^{(l)}) > 0$, the Jacobian of the representation with respect to the initial input features satisfies:
    \begin{equation}
    \left\| \frac{\partial \tilde{h}_v^{(L)}}{\partial x_v} \right\| \ge \prod_{k=1}^L \sigma_{\min}(W_{\text{skip}}^{(k)}) > 0
    \end{equation}
    guaranteeing that input feature sensitivity does not vanish across message-passing layers.
\end{enumerate}
\end{theorem}

\begin{proof}
We establish the proof in three parts: intrinsic feature erasure, gradient lower bounds under residual skip connections, and the specific application to identity multigraphs.

\paragraph{Part 1: Intrinsic Feature Erasure under Pure Neighborhood Aggregation.}
Let $v \in V$ be an arbitrary node in $G$ with initial feature vector $x_v = h_v^{(0)} \in \mathbb{R}^{d_0}$. Under the standard heterogeneous relational convolution layer:
\begin{equation}
h_v^{(1)} = \sigma \left( \sum_{r \in \mathcal{R}_{\text{in}}(\tau(v))} \bigoplus_{u \in \mathcal{N}_r(v)} \alpha_{vu}^{(1)} W_r^{(1)} h_u^{(0)} \right)
\end{equation}
Observe that the argument of $\sigma(\cdot)$ is a linear combination of neighbor states $\{h_u^{(0)} \mid u \in \mathcal{N}_r(v), r \in \mathcal{R}_{\text{in}}(\tau(v))\}$. In an access control multigraph lacking explicit self-loop relations on template nodes:
\begin{equation}
v \notin \mathcal{N}_r(v), \quad \forall r \in \mathcal{R}_{\text{in}}(\tau(v))
\end{equation}
Differentiating $h_v^{(1)}$ with respect to $x_v = h_v^{(0)}$:
\begin{equation}
\frac{\partial h_v^{(1)}}{\partial x_v} = \sigma'(\cdot) \cdot \sum_{r \in \mathcal{R}_{\text{in}}(\tau(v))} \sum_{u \in \mathcal{N}_r(v)} \alpha_{vu}^{(1)} W_r^{(1)} \frac{\partial h_u^{(0)}}{\partial h_v^{(0)}}
\end{equation}
Because distinct nodes have disjoint representations ($\partial h_u^{(0)} / \partial h_v^{(0)} = \mathbf{0}$ for all $u \neq v$), every term in the summation evaluates to zero:
\begin{equation}
\frac{\partial h_v^{(1)}}{\partial x_v} = \mathbf{0} \in \mathbb{R}^{d_1 \times d_0}
\end{equation}
Thus, at layer 1, the node's updated embedding $h_v^{(1)}$ is completely independent of its initial feature vector $x_v$. In certificate templates ($t \in V_{\text{Template}}$), this means that the 10 binary configuration flags $x_{\text{Template}}$ (which govern whether the template allows enrollee-supplied SANs, requires manager approval, or specifies client authentication EKUs) are completely discarded from $h_t^{(1)}$. 

Furthermore, if $v$ is a source-only entity satisfying $d_{\text{in}}(v) = 0$ (such as user accounts possessing outgoing enrollment rights but no incoming delegations), then $\mathcal{N}_r(v) = \emptyset$ for all $r$. By the identity of aggregation over an empty set ($\bigoplus(\emptyset) = \mathbf{0}$) and $\sigma(\mathbf{0}) = \mathbf{0}$, we obtain $h_v^{(l)} = \mathbf{0}$ for all $l \ge 1$.

\paragraph{Part 2: Gradient Lower Bounds via Parameterized Residual Skips.}
Now consider the layer formulation equipped with parameterized additive residual skip connections:
\begin{equation}
\tilde{h}_v^{(l)} = h_v^{(l)} + W_{\text{skip}}^{(l)} \tilde{h}_v^{(l-1)}
\end{equation}
where $\tilde{h}_v^{(0)} = x_v$ and $W_{\text{skip}}^{(l)} \in \mathbb{R}^{d_l \times d_{l-1}}$. Unrolling the recurrence from layer $L$ to layer $0$:
\begin{equation}
\tilde{h}_v^{(L)} = \left( \prod_{k=1}^L W_{\text{skip}}^{(k)} \right) x_v + \sum_{l=1}^L \left( \prod_{k=l+1}^L W_{\text{skip}}^{(k)} \right) h_v^{(l)}
\end{equation}
where by convention $\prod_{k=L+1}^L W_{\text{skip}}^{(k)} = I_{d_L}$. 

Taking the Jacobian of $\tilde{h}_v^{(L)}$ with respect to $x_v$:
\begin{equation}
\frac{\partial \tilde{h}_v^{(L)}}{\partial x_v} = \prod_{k=1}^L W_{\text{skip}}^{(k)} + \mathcal{J}_{\text{graph}}(v)
\end{equation}
where $\mathcal{J}_{\text{graph}}(v) = \sum_{l=1}^L \left( \prod_{k=l+1}^L W_{\text{skip}}^{(k)} \right) \frac{\partial h_v^{(l)}}{\partial x_v}$ represents the indirect gradient flowing through cyclical graph paths (e.g., $v \to u \to v$).

By the reverse triangle inequality for operator norms:
\begin{equation}
\left\| \frac{\partial \tilde{h}_v^{(L)}}{\partial x_v} \right\| \ge \left\| \prod_{k=1}^L W_{\text{skip}}^{(k)} \right\| - \|\mathcal{J}_{\text{graph}}(v)\|
\end{equation}
For the direct skip path, the minimum singular value of the matrix product satisfies:
\begin{equation}
\sigma_{\min} \left( \prod_{k=1}^L W_{\text{skip}}^{(k)} \right) \ge \prod_{k=1}^L \sigma_{\min}(W_{\text{skip}}^{(k)})
\end{equation}
Assuming full-rank initialization with $\sigma_{\min}(W_{\text{skip}}^{(k)}) > 0$ for all $k \in \{1, \dots, L\}$ (as standard in linear projection layers initialized via Xavier or He uniform initialization), the skip connection establishes a guaranteed, non-vanishing path preserving initial attributes. Following the analytical framework of GCNII (Chen et al., 2020)~\cite{chen2020simple}, this prevents over-smoothing and feature decay, ensuring that:
\begin{equation}
\left\| \frac{\partial \tilde{h}_v^{(L)}}{\partial x_v} \right\| \ge \prod_{k=1}^L \sigma_{\min}(W_{\text{skip}}^{(k)}) > 0
\end{equation}

\paragraph{Part 3: Implication for Active Directory Certificate Template Classification.}
In Active Directory Certificate Services, certificate template classification is a joint decision requiring both:
\begin{enumerate}
    \item Intrinsic configuration flags $x_t \in \mathbb{R}^{10}$ (e.g., whether \texttt{CT\_FLAG\_ENROLLEE\_SUPPLIES\_SUBJECT} and Client Authentication EKUs are enabled).
    \item Topological reachability from unprivileged principals ($\exists u \in V_{\text{low-priv}} \rightsquigarrow t$).
\end{enumerate}
Without residual skip connections ($W_{\text{skip}} = 0$), $h_t^{(1)}$ completely erases $x_t$. While $x_t$ can theoretically influence $h_t^{(2)}$ through 2-hop cycles ($t \to u \to t$), this signal is severely diluted by attention normalization across all adjacent principals and relations. Consequently, the network loses direct access to the template's configuration flags, explaining the empirical collapse from Macro-F1 = $0.9986$ to $0.4768$ observed in Section~\ref{sec:ablation_results}. Parameterized residual skip connections guarantee that the 10 binary flags are directly retained in the final node embedding, preserving classification efficacy.
\end{proof}

\begin{corollary}[Necessity in Active Directory Identity Graphs]
In enterprise identity graphs, where certificate templates rely heavily on intrinsic configuration attributes and user principals frequently act as low in-degree sources in administrative authorization subgraphs, parameterized residual skip connections are mathematically required to prevent feature decay and gradient vanishing across multi-hop relational convolutions.
\end{corollary}
